\section{Limitations and Next Steps}

\subsection{Current Limitations}
\begin{itemize}
  \item Most dashboard interpretations still come from single-feature EDA, so
  the findings describe strong segment signals rather than full causal
  explanations.
  \item The model metrics require careful reading because the churn class is
  imbalanced; accuracy is less informative than recall, precision, F1-score,
  ROC AUC, and PR-AUC for this task.
  \item Several numeric variables are grouped for dashboard readability, which
  improves interpretation but can reduce predictive detail for the model.
  \item Some likely churn drivers are not available in the dataset, including
  detailed payment history, plan-change events, discount exposure, competitor
  offers, and support-ticket content.
  \item The final dashboard prioritizes the course techniques that best support
  this churn story. More optional views, such as small multiples or additional
  distribution plots, remain possible future extensions.
  \item Some visual polish remains possible, especially typography, page-level
  titles, and a final contrast audit for every chart.
\end{itemize}

\subsection{Next Steps}
\begin{enumerate}
  \item Re-check all churn-rate values against the latest cleaned dataset.
  \item Add more multivariate analysis after the single-feature EDA is stable.
  \item Review Random Forest feature importance and high-risk customer examples.
  \item Test raw numeric features together with grouped features and compare
  Random Forest with XGBoost, LightGBM, or CatBoost.
  \item Add PR-AUC, threshold analysis, and class-imbalance handling to the
  modeling evaluation if the predictive workflow is extended.
  \item If more time is available, refine chart titles, verify color contrast,
  and audit churn-rate denominators directly in the dashboard artifact.
  \item Add more final insight paragraphs after the group completes the oral
  presentation narrative.
\end{enumerate}

\section{Conclusion}
This report connects the churn dataset, EDA workflow, feature grouping, Random
Forest outputs, and final Power BI dashboard into one analysis narrative. The
main interpretation separates distribution from risk: customer counts explain
segment size, while churn rates and average predicted probabilities explain
which groups need attention. The final dashboard applies course concepts from
visual encoding, graphical perception, table-data visualization, figure design,
interaction, and storytelling. In the current results, customer-experience
features such as satisfaction, complaints, service calls, late payments, and
contract type provide the clearest signals for retention prioritization, while
the model supports ranking customers rather than replacing business judgement.
